\documentclass{report}
\usepackage{amsfonts, amsmath}
\newcommand\R{{\mathbb R}}
\newcommand\Q{{\mathbb Q}}
\newcommand\N{{\mathbb N}}
\pagestyle{empty}
\begin{document}
\large
\centerline{\Huge \bf Analisi Matematica I}
\vskip.2cm
\centerline{\Huge Primo Appello}
\renewcommand{\thefootnote}{ } 
\footnote{\sl Avete 3 ore di tempo. Ogni esercizio
vale sei punti. La sufficienza si ottiene con un punteggio $\geq$
18. {\bf Solo le risposte chiaramente giustificate saranno prese in
considerazione.} {\scshape Elaborati illeggibili o confusi verranno ignorati.}
{\sffamily La
sintesi sar\`a particolarmente apprezzata.}}
\renewcommand{\thefootnote}{\arabic{footnote}} 

\vskip1cm
\begin{enumerate}
\item Si trovino gli $x\in \R$ tali che
\[
|2x-1|-x<1.
\]
\item Si disegni l'insieme $C=A\cap B$ dove $A=\{(x,y)\in\R^2\;|\;
x^2+y^2\leq 1\}$ e $B=\{(x,y)\in\R^2\;|\;y\geq x\}$. Dire se il punto
$(-1/2,0)$ appartiene o non appartiene a $C$.
\item Si disegni il grafico della funzione $f:\R\to\R$ definita da
\[
f(x)=\frac{e^{x}+e^{-x}}2.
\]
Si trovi il numero dell soluzioni dell'equazione
\[
f(x)-1-x=0.
\]
\item Sia $\ell_1$ il segmento rettilineo che congiunge i punti
$(0,1)$ e $(x,0)$, ed $\ell_2$ quello congiungente i punti $(x,0)$ e
$(2,1)$. Si determinino i valori di $x\in[0,2]$ per cui la somma delle
lunghezze di $\ell_1$ e $\ell_2$ \`e minima.
\item Si studi il seguente limite
\[
\lim_{n\to\infty}\frac{n\sin\frac 1n -1}{n^2}.
\]
\item Si scriva una funzione $f:\R\to\R$, derivabile in tutti i punti,
tale che $f(x)=0$ per tutti gli $x<0$ e $f(x)=x$ per tutti gli $x>1$. 
\end{enumerate}
\end{document}
